\documentclass[11pt]{article}

\usepackage[a4paper,margin=1in]{geometry}
\usepackage{amsmath}
\usepackage{booktabs}
\usepackage{tabularx}
\usepackage[hidelinks]{hyperref}

\title{Current Numerical Results for the Bermudan Put Experiments}
\date{13 March 2026}

\begin{document}
\maketitle

\section*{Purpose of this note}

This note collects the numerical results currently available from the standalone scripts.
The idea is to have one short working document that can be shared and discussed easily.
At this stage, the main point is to show where the benchmark stands, how the generalized
double mean-reverting (gDMR) setup is behaving, and what changed after the recent tuning.

\section*{1. Paper benchmark: Heston Bermudan put}

The \texttt{paper\_setup} folder is kept as a benchmark against the Heston Bermudan put case
reported by Farahany, Jackson, and Jaimungal. In this case we do have a published finite-difference
reference value, so the comparison is quite direct.

\begin{table}[h]
\centering
\caption{Benchmark comparison for the paper-style Heston setup.}
\begin{tabular}{lrrrr}
\toprule
Estimator & Published value & Current run & $|$Run $-$ FD$|$ & Run $-$ published \\
\midrule
LSMC direct   & 1.4494 & 1.458896 & 0.008196 & +0.009496 \\
LSMC low      & 1.4487 & 1.452544 & 0.001844 & +0.003844 \\
Hybrid direct & 1.4530 & 1.452988 & 0.002288 & -0.000012 \\
Hybrid low    & 1.4529 & 1.445493 & 0.005207 & -0.007407 \\
\bottomrule
\end{tabular}
\end{table}

The finite-difference reference for this benchmark is 1.4507. The most encouraging point here is
that the hybrid direct estimator is essentially on top of the published hybrid direct value, while
the other estimators remain in the same numerical range as the paper.

\section*{2. Current gDMR baseline: $\delta_1=\delta_2=0.50$}

For the gDMR case there is no published finite-difference benchmark available in the project, so the
current practical reference is the large LSMC direct run. The default baseline uses the following setup.

\begin{table}[h]
\centering
\caption{Main numerical settings for the current gDMR baseline run.}
\begin{tabular}{lr}
\toprule
Quantity & Value \\
\midrule
$S_0$ & 100 \\
$K$ & 100 \\
$T$ & 1 \\
$r$ & 0.03 \\
$v_0$ & 0.04 \\
$v'_0$ & 0.04 \\
$\kappa_1$ & 2.0 \\
$\kappa_2$ & 1.0 \\
$\theta$ & 0.04 \\
$\xi_1$ & 0.35 \\
$\xi_2$ & 0.20 \\
$\delta_1$ & 0.50 \\
$\delta_2$ & 0.50 \\
LSMC paths & 1,000,000 \\
Exercise dates & 100 \\
Euler steps & 600 \\
Hybrid paths $N$ & 30,000 \\
Hybrid low paths $N_{\mathrm{low}}$ & 30,000 \\
Asset grid points $N_S$ & 181 \\
Hermite nodes & 64 \\
\bottomrule
\end{tabular}
\end{table}

\begin{table}[h]
\centering
\caption{Final baseline results for the current gDMR setup.}
\begin{tabular}{lrrr}
\toprule
Estimator & Price & Standard error & Relative error \\
\midrule
LSMC direct   & 6.384535 & 0.007943 & reference \\
LSMC low      & 6.387249 & 0.007939 & 0.04\% vs LSMC direct \\
Hybrid direct & 6.421758 & 0.000652 & 0.58\% vs LSMC direct \\
Hybrid low    & 6.265070 & 0.045711 & 1.87\% vs LSMC direct \\
\bottomrule
\end{tabular}
\end{table}

The main result here is that the hybrid direct estimator is now within 1\% of the LSMC direct reference.
That is a clear improvement over the earlier hybrid runs and is, in my view, the strongest current numerical
message from the gDMR experiments.

\begin{table}[h]
\centering
\caption{Tuning path for the hybrid method in the gDMR baseline.}
\begin{tabularx}{\textwidth}{lXrrr}
\toprule
Run & Key settings & Hybrid direct & Hybrid low & Direct rel. error vs LSMC direct \\
\midrule
Tier 1 & 10,000 paths, 121 grid points, 48 nodes, wide asset grid & 7.335311 & 6.621170 & 14.89\% \\
Tier 2 & 20,000 paths, 161 grid points, 48 nodes, wide asset grid & 6.987310 & 6.511937 & 9.44\% \\
Tier 3 & 40,000 paths, 201 grid points, 64 nodes, wide asset grid & 6.804899 & 6.424076 & 6.58\% \\
Truncated Tier 2 & 20,000 paths, 161 grid points, 48 nodes, truncation added, wide asset grid & 6.756928 & 6.281426 & 5.84\% \\
Previous default & 20,000 paths, 161 grid points, 48 nodes, truncation and tighter asset grid & 6.486941 & 6.279662 & 1.60\% \\
Current default & 30,000 paths, 181 grid points, 64 nodes, truncation and tighter asset grid & 6.421758 & 6.265070 & 0.58\% \\
\bottomrule
\end{tabularx}
\end{table}

This progression suggests that the most important improvements were not cosmetic. The gains came mainly
from three choices: a stronger Monte Carlo reference, truncation of the volatility regression domain, and
a tighter asset grid around the strike.

\section*{3. Additional gDMR experiment: $\delta_1=\delta_2=0.75$}

I also ran a second gDMR experiment with $\delta_1=\delta_2=0.75$, keeping the same tuned numerical
configuration as in the current baseline. This gives a simple way to see how the price level changes when
both elasticity parameters are increased.

\begin{table}[h]
\centering
\caption{Results for the equal-delta case $\delta_1=\delta_2=0.75$.}
\begin{tabular}{lrrr}
\toprule
Estimator & Price & Standard error & Relative error \\
\midrule
LSMC direct   & 6.653560 & 0.007822 & reference \\
LSMC low      & 6.655062 & 0.007819 & 0.02\% vs LSMC direct \\
Hybrid direct & 6.696152 & 0.000375 & 0.64\% vs LSMC direct \\
Hybrid low    & 6.545257 & 0.045080 & 1.63\% vs LSMC direct \\
\bottomrule
\end{tabular}
\end{table}

The 1\% direct-error target is still met in this equal-delta case. So, at least for the current parameter set,
the hybrid direct estimator remains stable after moving from $\delta=0.50$ to $\delta=0.75$.

\begin{table}[h]
\centering
\caption{Comparison of the two gDMR cases currently tested.}
\begin{tabular}{lrrr}
\toprule
Estimator & $\delta=0.50$ & $\delta=0.75$ & Change \\
\midrule
LSMC direct   & 6.384535 & 6.653560 & +0.269025 \\
LSMC low      & 6.387249 & 6.655062 & +0.267813 \\
Hybrid direct & 6.421758 & 6.696152 & +0.274394 \\
Hybrid low    & 6.265070 & 6.545257 & +0.280187 \\
\bottomrule
\end{tabular}
\end{table}

Numerically, all four estimators move upward when the model is changed from $\delta_1=\delta_2=0.50$ to
$\delta_1=\delta_2=0.75$. At the same time, the agreement between the direct estimators remains quite good.

\section*{4. Main points to discuss}

\begin{itemize}
\item The paper benchmark looks healthy: the standalone implementation reproduces the published Heston numbers reasonably well, especially for the hybrid direct estimator.
\item In the gDMR case, the strongest current result is that the tuned hybrid direct estimator now sits within 0.58\% of the LSMC direct reference.
\item The low estimator is still visibly weaker than the direct estimator in the gDMR experiments, so this remains a natural place for further improvement.
\item The equal-delta experiment with $\delta_1=\delta_2=0.75$ raises the price level, but it does not break the current numerical behavior of the direct hybrid method.
\end{itemize}

\section*{5. Short conclusion}

At the moment, I think the main story is quite clear. The benchmark case gives confidence that the implementation
is behaving sensibly, and the tuned gDMR setup now produces a hybrid direct value that is close to the Monte Carlo
reference. That feels like a solid basis for discussion with the supervisor, even if the low estimator still needs
more work.

\end{document}
